% SPDX-FileCopyrightText: 2026 Will Cook
% SPDX-License-Identifier: CC-BY-4.0
\section{Learning to explain a particular argument}\label{sec:reading}
\papersectiontarget{exposition-reading}

Suppose a proof constructs a sequence by repeatedly choosing a digit. A sentence
saying that the choice can be repeated explains how to obtain every finite
prefix. It may leave the reader to supply the reason why the resulting infinite
series has the prescribed sum. Replacing the sentence by smoother prose does
not address that difficulty. The explanation needs the invariant for the
remainder and the fact that its bound tends to zero.

This example, developed in Section~\ref{sec:remainder}, suggests where to look
for a useful model. Another paper that constructs an infinite sum through
controlled finite remainders can show how to separate the choice from the
limiting argument. It need not study the same sequence. Its relevance lies in
an identifiable part of the proof, and the comparison ends where the arguments
differ. The model contributes an expository choice; the new paper must supply
its own mathematics.

General advice remains useful. Halmos discusses audience, organization and the
selection of examples~\cite[\S\S3--4]{halmos}. Knuth, Larrabee and Roberts discuss
the introduction of symbols, the connection between sentences and the reader's
need to know why a step is being taken~\cite[\S1]{knuth}. Gowers examines the
placement of examples relative to unfamiliar abstractions~\cite{gowers}, while
Tao cautions against sacrificing usefulness to excessive optimization of a
paper~\cite{tao}. Such advice identifies questions to ask. A nearby mathematical
argument shows how an author has answered them in a particular setting.

The method described here grew out of revisions to short papers and longer
research records for eight Erd\H{o}s problems. Most research and drafting in
this project was performed by AI agents. Some agents worked from frozen packets;
others had access to the live repository and integrated proposed changes.
Their judgements are editorial evidence, with the limitations described below.
They are not observations of how an independent mathematical audience read the
papers. The examples support an inspectable procedure for revision rather than
a claim that the procedure has already been shown to produce better readers or
better mathematics.

\subsection{Choose sources by the work the prose must do}

Before collecting stylistic examples, identify the manuscript's subject, genre
and reader. An irrationality note may need a compact route from an assumed
rational sum to a positive integer smaller than one. A geometric paper may need
to explain which component is being measured and why a map preserves the
quantity under discussion. A survey has to distinguish several viewpoints; a
short theorem paper may need to develop just one. These differences affect
terminology, notation and the amount of preparation a proof needs.

Choose a small group of human-authored papers that addresses these needs.
Include a paper close in subject for vocabulary and attribution, and, when
useful, one close in proof structure for the explanation of a construction or
estimate. A famous paper from a remote field can offer a comparison, but it
cannot settle the conventions of the target field. Nor does frequent usage
make a term appropriate when its definition differs from the object at hand.

Read the original passage in a named version. An abstract can establish the
name used for an object; it seldom reveals the pacing of the proof. For that,
read the argument around the sentence, including the statement it proves and
the hypotheses it invokes. Record the section, theorem or equation label and
page. Preprint versions can have different numbering from a published article,
and a source archive may differ from the accompanying PDF. The file's digest
identifies the supplied bytes; a public link tells the reader where to seek
the work. Neither should silently stand in for the other.

The record should distinguish three acts. A packet may contain a source. Its
recipient may declare that particular pages were read. An integrating reviewer
may later inspect the quoted passage. These are different facts. Later
inspection can verify a locator or expose a poor analogy; it cannot establish
what the earlier recipient actually read.

\subsection{Read sentences as mathematical relations}

A useful reading pass asks what each choice allows the reader to infer.
Consider the placement of a hypothesis. A theorem statement should quantify
the objects and state the assumptions needed for its conclusion. Inside the
proof, recalling one of those assumptions can explain a step that would
otherwise look automatic. A phrase such as ``by positivity'' helps only if
the reader can identify the positive quantity and the inference that uses it.

Connective words deserve the same attention. ``Since'' introduces a reason;
``hence'' asserts that the conclusion follows; ``provided that'' makes a
condition visible; ``it remains to prove'' identifies unfinished work. These
words are not interchangeable devices for varying a paragraph. A revision
that changes ``provided that'' to ``hence'' can change the logical force of
the sentence while leaving every displayed formula intact.

Parallel syntax helps compare parallel claims. A subordinate clause can keep
a dependency beside the assertion that uses it. There is no preferred ratio
of short to long sentences: a short sentence may hide its antecedent, while
a longer one can keep a condition attached to its consequence. The point of
studying an author's construction is to understand that relation, then write
original prose appropriate to the new argument.

Notation should receive this close reading as well. Ask when a symbol first
becomes necessary, what repeated work it saves and which nearby quantities
the reader might confuse with it. A new name for a standard object creates an
additional translation. Conversely, removing an established term such as
``Stieltjes moment sequence'' can make a manuscript less connected to its
literature~\cite{wangzhu}. A vocabulary pass should compare definitions and
uses, rather than delete unfamiliar words merely because they are technical.

Proof pacing is the allocation of explanation to mathematical difficulty.
The reader may need to see the accuracy required of an estimate before its
derivation, or the obstruction that motivates a parameter choice. Several
lines of routine substitution can be shorter than one indispensable sentence
about why a limit exists. The manuscript should give the latter its space.
These choices can be examined in actual passages, without imitating an
author's voice or turning a preferred cadence into a rule.
